\section{The actual theta-index normalization and the mixed-character gate}
\label{sec:theta-cohomology-interface}

The regular construction of AC-G27 can be compared with the
literal arithmetic normalization in the supplied source.  This
comparison gives a bounded lattice loss, but the cohomology
bound remains a separate input.  Throughout this section the
source version is the immutable commit
\texttt{7cbf7a0c32fca1badf05931879a37c3f8a2b03b4}.
The newly recovered files and their hashes are recorded in the
companion notes and acquisition manifest.

\begin{theorem}[AC-G28: cubic projectors in the source index units]
\label{AC-G28}
Let $R$ be the ring of integers of a finite extension of
$\mathbb Q_\ell$ containing $\omega$, let $V$ be a finite-dimensional
vector space over its fraction field, and let $M\subset V$ be a
full $R$-lattice.  Suppose $U(M)=M$ and $U^3=1$.
For $J\subset\{0,1,2\}$ put
\[
 P_J=\sum_{j\in J}\frac13\sum_{a=0}^2\omega^{-ja}U^a,
 \qquad V_J=P_JV,
 \qquad M_J^+=P_JM,\quad M_J^-=M\cap V_J.
                                                        \tag{G28.1}
\]
Then
\[
 3M_J^+\subset M_J^-\subset M_J^+,
 \qquad
 M_J^+=V_J\cap(M+V_{J^c}).                              \tag{G28.2}
\]
Writing $r_J=\dim V_J$, the normalized length used in
\texttt{DataN.lenUnit} satisfies
\[
 0\le\frac{\operatorname{length}_R(M_J^+/M_J^-)}
      {\operatorname{length}_R(R/\ell R)}
 \le\begin{cases}
 r_J,&\ell=3,\\
 0,&\ell\ne3.
 \end{cases}                                            \tag{G28.3}
\]
The same inclusion bound holds after intersecting with a
rational subspace or passing to a rational quotient.  If the
integral $C_3$-action is retained through tensor, Weyl, dual,
or other functorial integral constructions and the character
projection is performed only at the end, (G28.2)--(G28.3)
still hold with the constant $3$, independently of their degree.
In contrast, applying a homogeneous polynomial functor of degree
$d$ to the two lattices in (G28.2), after projection, gives the
general inclusion $3^d F(M_J^+)\subset F(M_J^-)$.

If $M=R^a\oplus R[C_3]^b$, with trivial action on $R^a$, then
for every nonempty proper $J$ one has the sharper exact formula
\[
 M_J^+/M_J^-\simeq(R/3R)^b,
 \qquad \operatorname{lenUnit}_3(M_J^+/M_J^-)=b
                      \quad(\ell=3).                  \tag{G28.4}
\]
The theta permutation lattice of AC-G27 therefore has this
normalized length $b=(r^3-r)/3$, where $r=d^g$, for every
nontrivial proper character selection.  This projected-to-intersected
length is a different invariant from the full Fourier determinant
index of AC-G27.
\end{theorem}
\begin{proof}
The polynomial $X^3-1$ is separable over the fraction field.
The three displayed character operators are mutually orthogonal
idempotents and sum to the identity.  Thus $V=V_J\oplus V_{J^c}$.
Every coefficient of $3P_J$ belongs to $R$, so $3P_JM\subset M$;
it also lies in $V_J$.  Conversely $P_J$ is the identity on
$M\cap V_J$.  This proves the inclusions.

If $x=P_Jm$, then $x=m-P_{J^c}m\in V_J\cap(M+V_{J^c})$.
If $x=m+y$ belongs to that intersection, with $m\in M$ and
$y\in V_{J^c}$, applying $P_J$ gives $x=P_Jm$.
This proves the second equality in (G28.2), exactly the
projected/intersected shape in \texttt{NData.MtUp} and
\texttt{NData.MtLow}.

Both $M_J^+$ and $M_J^-$ are free of rank $r_J$.
Their quotient is a quotient of $M_J^+/3M_J^+$, so its
length is at most $r_Jv_R(3)$.  The denominator in
\texttt{lenUnit} is $v_R(\ell)$.  At $\ell=3$ these
valuations agree, and at other primes $3$ is a unit.
This proves (G28.3), including finite length; no manipulation
of the source's $\mathrm{toNat}$ at infinite length is involved.

For any rational subspace $W$, multiplying an element of
$M_J^+\cap W$ by $3$ places it in $M_J^-\cap W$.
Applying a linear quotient map gives the corresponding statement
for image lattices.  For any newly constructed integral module
with the induced order-three action, the same projector proof
starts afresh and has the same denominator $3$.  No estimate of
its dimension or tensor degree enters this argument.
If instead $F$ is a homogeneous polynomial functor of degree
$d$ applied to an already projected inclusion, the maps
$3\operatorname{id}:M_J^+\to M_J^-$ and
$\operatorname{id}:M_J^-\to M_J^+$ induce maps whose composite
is $F(3\operatorname{id})=3^d\operatorname{id}$.
Interpreting their images in the common rational representation
gives the asserted degree-dependent inclusion.  The latter
claim is restricted to such polynomial functors; the
degree-independent projector assertion only requires an
integral order-three action.

It remains to compute one regular orbit.  Write its basis as
$e_0,e_1,e_2$, with $Ue_a=e_{a+1}$, the suffixes taken modulo
$3$.  For a singleton $J=\{j\}$, an eigenvector is
$v_j=\sum_a\omega^{-ja}e_a$.  Its coefficients
are units, so $M\cap V_J=Rv_j$, whereas $P_JM=(1/3)Rv_j$.
The quotient is $R/3R$.
For a two-element $J$, let $k$ be its complement.  The
linear functional defining $V_J$ has unit coefficients and
takes the value $3$ on a generator of the complementary
eigenline.  It follows that
\[
 M\big/\big((M\cap V_J)+(M\cap V_{\{k\}})\big)
              \simeq R/3R.
\]
Projection identifies this quotient with $M_J^+/M_J^-$.
The trivial summands contribute nothing, and direct sums give
(G28.4).  The integral orbit decomposition of AC-G27 now gives
the theta-lattice assertion.
\end{proof}

\begin{corollary}[The exact height and global index normalization]
\label{AC-G28-normalization}
Let $F$ be a totally imaginary number field of degree $D$,
with the same Hermitian metrics on two full $\mathcal O_F$-lattices
$L\subset M$ of rank $n>0$.  Suppose
$3^aM\subset L\subset M$, for an integer $a\ge0$.
In the convention of
\texttt{Formalism.HarderNarasimhan.StatementPl.size} and
\texttt{muDef},
\[
 0\le\mu(M)-\mu(L)
   =\frac{2}{Dn}\log[M:L]\le2a\log3.                  \tag{G28.5}
\]
The same bounds hold for the differences of the maximum
slopes and of the minimum slopes.  Thus a final cubic
character projection costs at most $2\log3$ in each of
these slope conventions.

There is also a direct conditional comparison in the precise
\texttt{LinkN\_at} units.  Suppose an actual source row has
\texttt{etaJ\_tie} identifying its local indices with comparisons
as in (G28.1), and suppose for every coefficient-field place
$w\mid3$ the sum of their selected ranks is at most
$d_0=\texttt{dZ0}\,j$.  The additional index contribution is
then at most
\[
 \sum_{w\mid3}\frac{\log Nw}{[F:\mathbb Q]}
                 e(w/3)\,\eta_J(j,w)
              \le d_0\log3.                            \tag{G28.6}
\]
It is zero at all other coefficient primes.  Hence the
additional term inside the rank-normalized bracket of
\texttt{LinkN\_at} is at most $\log3$, once that row dictionary
and rank comparison have been proved.
\end{corollary}
\begin{proof}
The two covolumes differ by the additive-group index
$[M:L]$.  All degree, discriminant, and rank factors in the
definition of \texttt{ht} are the same.  Since $F$ is totally
imaginary, $2\,\texttt{nPl}=D$; the source definitions give
\[
 \operatorname{size}(M)=\operatorname{ht}(M)^{1/(Dn)},
 \qquad \mu(M)=-2\log\operatorname{size}(M).
\]
Thus the middle equality in (G28.5) follows.  The quotient
$M/L$ is a quotient of $M/3^aM$, whose cardinality is
$3^{aDn}$, proving the inequalities.

Intersecting with any rational subspace preserves the
inclusions and the constant $3^a$.  Passing to the image in
any rational quotient also preserves them.  Apply (G28.5)
to each such intersection or quotient.  The maximum slope
is the supremum of the intersection slopes: a sublattice
with a fixed rational span has no larger slope than the
intersection with that span.  Similarly the minimum slope
is the infimum of the saturated quotient slopes.  Taking
these suprema and infima proves the corresponding bounds.

For (G28.6), (G28.3) and the assumed rank comparison imply
$\eta_J(j,w)\le d_0$.  All summands are nonnegative.
Writing $Nw=3^{f(w/3)}$ and using
$\sum_{w\mid3}e(w/3)f(w/3)=D$ yields the stated bound.
The source \texttt{etaSum} uses the weight
$\texttt{lwP}(w)\,\texttt{ramIdx}(w)$, with no factor $2$;
\texttt{DataArith.sizeOfVal} uses twice that weight.
This accounts for $\log3$ in (G28.6) and $2\log3$ in
(G28.5).  Neither is a normalization of a spectral trace.
\end{proof}

\begin{remark}[What the theta ties do and do not identify]
The exact source locators for the preceding comparison are
\texttt{DataN.lean}: \texttt{lenUnit}, \texttt{NData.MtUp},
\texttt{MtLow}, \texttt{etaLoc};
\texttt{NStageTrivialCharacterCoeffEq.lean}: the field
\texttt{NStage.etaJ\_tie}; and
\texttt{DataParams\_2}: \texttt{SizeData\_2.dZ0},
\texttt{lwP}, \texttt{sz}.  The last two displayed filenames
with an ASCII suffix refer to the supplied files whose actual
suffix is the Unicode subscript $2$.
The source \texttt{DataTheta.ThetaX.TieRowIsom} instead says
that a Petersson product scaled by \texttt{mY} equals the
given row inner product; \texttt{TieMY} expresses
\texttt{mY} as a product of explicit parabolic indices.
The source \texttt{DataCoeff.normFock} divides a Fock pairing
by the pairing of the chosen highest-weight vector with itself.
Neither of these declarations divides a determinant index
by the number of theta sections.
The inequality (G28.6) is therefore a conditional comparison
with the actual row index: the identification of the new
regular construction with those rows, including the selected
rank bound, has not been supplied by AC-G27 or by the
archimedean orthogonality of its projectors.
\end{remark}

\begin{theorem}[AC-G29: scalar cancellation and the actual cohomology gate]
\label{AC-G29}
Let $G$ be a group, $R$ a commutative ring, and
$\kappa:G\to R^\times$ a character.  For $R$-representations
$V,W$, tensoring both with the free rank-one character
$R(\kappa)$ gives an integral, $G$-equivariant identification
\[
 \operatorname{Hom}_R(V\otimes R(\kappa),
                     W\otimes R(\kappa))
       \simeq\operatorname{Hom}_R(V,W).                 \tag{G29.1}
\]
In particular, the literal family in \texttt{NData.chi\_def},
\[
          \chi_c=\chi_1\varepsilon^{-(c-1)},
                         \qquad 1\le c\le t,
\]
has pairwise coefficients
\[
 \operatorname{Hom}(R(\chi_a),R(\chi_b))
                =R(\varepsilon^{a-b}).                 \tag{G29.2}
\]
A common cubic twist of this family leaves these coefficients
unchanged, including at the coefficient prime $3$.

For comparison, put $P=R[C_3]$ with its regular action.
Under conjugation by the regular permutation,
\[
              \operatorname{End}_R(P)\simeq P^{\oplus3}
                                                        \tag{G29.3}
\]
integrally.  After extending scalars to a characteristic-zero
splitting field, the mixed family
$R(\chi_1\kappa^i\varepsilon^{-(a-1)})$, $i=0,1,2$,
has cross coefficients
\[
       R(\kappa^{j-i}\varepsilon^{a-b}).                \tag{G29.4}
\]
Thus retaining regular cubic blocks produces the required
mixed cubic coefficients, whereas a common scalar twist
cannot produce them through (G29.2).

In the supplied source the existing arithmetic tie is
\texttt{NData.sha\_eq}: at each finite length and
$0<e<t$, it identifies \texttt{Rows.shaLen} with the
strict cohomology length of
$\mathbb Q_\ell/\mathbb Z_\ell(-e)$ over the actual field
\texttt{Host.E}.  It is not a tie to (G29.4).
The source \texttt{LinkN\_at} has the direction
\[
 \eta\le d_0\left(C\,\mathcal E+
           \sum_\ell\log\ell\sum_{0<e<t}
                    \operatorname{shaLen}_\ell(e)\right),
                                                        \tag{G29.5}
\]
with an explicitly defined geometric quantity $\mathcal E$.
Consequently an upper bound on the additional index from
AC-G28 supplies no upper bound on the cohomology sum in
(G29.5).  That sum is bounded separately by
\texttt{LinkV\_at} or by the varying-field \texttt{VVle}
input to \texttt{VAt}/\texttt{VWinAt}.
\end{theorem}
\begin{proof}
Choose the standard generator of the free rank-one module
$R(\kappa)$.  It identifies the underlying Hom modules in
(G29.1).  For a map $f$, the Galois action on the twisted Hom
is
\[
 g\cdot f=(\kappa(g)\rho_W(g))\,f\,
                 (\kappa(g)\rho_V(g))^{-1}
            =\rho_W(g)f\rho_V(g)^{-1}.
\]
This proves equivariance without dividing by the order of
the character.  Inverting the tensor functor by
$R(\kappa^{-1})$ also shows that it is exact and preserves
extension classes and all restriction maps.  For rank-one
characters the Hom character is $\chi_b\chi_a^{-1}$,
which gives (G29.2).  When $b>a$, its negative Tate weight
is precisely $e=b-a$ in the source \texttt{qzTw} convention.

For (G29.3), index a basis of $P$ by $i\in\mathbb Z/3$.
The matrix units $E_{ij}$ form a basis of its endomorphism
module.  Conjugation takes $E_{ij}$ to $E_{i+1,j+1}$.
There are exactly three free orbits, distinguished by
$j-i$.  Each is an integral regular module, proving (G29.3)
over every $R$, including residue characteristic $3$.
The rank-one ratio calculation gives (G29.4).
For finite projective $V,W$, one can equivalently identify
the regular-block Hom with
\[
 \operatorname{Hom}_R(V\otimes P,W\otimes P)
       \simeq\operatorname{Hom}_R(V,W)
                    \otimes\operatorname{End}_R(P).
\]
This expresses the new mixed coefficient before any
nonintegral character splitting.

The asserted source identification is read directly from
\path{GaloisCohomology.Statement.LData}: its finite set
$S$ is required to equal the primes above $\ell$, and its
character must pull back to the cyclotomic character.
The definitions are
$\texttt{qzTw}(e)=\mathbb Q_\ell/\mathbb Z_\ell(-e)$,
$\texttt{Sha}=\texttt{Sha1}$ for the local family at $S$,
and $\texttt{shaLen}=\operatorname{length}_{\mathbb Z_\ell}
\texttt{Sha}$.
In \texttt{NData}, \texttt{ldat} is such a datum for
$\texttt{H.E}$; \texttt{sha\_eq} gives the equality to the
row field whenever the actual length equals a natural number.
This ties the row variable to arithmetic when finite.  The
abstract \texttt{Rows.Data} countermodel of AC-G25 therefore
applies to that raw structure, and is not a countermodel to
a fully tied finite \texttt{NData}.

Finally \texttt{GlobalTheoremInputs.LinkN\_at} is exactly
(G29.5), under summability of its \texttt{vvTerm}.
Even numerically, an inequality
$0\le\eta\le d_0(E+V)$ with $d_0>0$ and $E,V\ge0$
cannot imply an upper bound on $V$: for example fix
$\eta=0$, $d_0=1$, $E=0$, and let $V$ be arbitrarily large.
This is a counterexample to reversing that inequality,
not an assertion that actual cohomology lengths in a fixed
number field can be assigned arbitrarily.
The independently declared \texttt{LinkV\_at} bounds the
cohomology sum from above.  The definitions
\texttt{ExistsPageN.VVle}, \texttt{VAt}, and
\texttt{VWinAt} quantify over the source host fields and
their windows.  Neither (G28.3) nor (G29.1) changes those
quantifiers or proves the required bound.
\end{proof}

\begin{proposition}[Exact transport of the local Hecke fragment]
\label{AC-G29-Hecke}
Let $R$ be the integers of a finite extension of $\mathbb Q_3$
containing $\omega$, and write $e_R=v_R(3)$.
Let $T\subset R^X$ be an $R$-subalgebra, where $X$ is finite
and nonempty, such that $T$ is local and separates the points
of $X$.  Let $M$ be a finite free $R$-module with a faithful
$T$-action.  The regular cubic extension
\[
 T'=T[C_3]\subset R^{X\times\{0,1,2\}},
             \qquad M'=M\otimes_R R[C_3]               \tag{G29.6}
\]
again gives a local algebra separating its rational
eigencharacters, acting faithfully on a finite free module.
These are the corresponding local algebra clauses of
\texttt{NData}; no global Hecke or automorphic realization is
asserted here.

Let $P$ be the rational projector associated with any subset
of $X$, and put
\[
 A=P(M),\qquad B=M\cap\operatorname{im}P,
                 \qquad r=\dim\operatorname{im}P.
\]
Select one cubic character in (G29.6).  For the joint
projector $P\otimes P_j$ the new projected and intersected
lattices are, respectively,
\[
           (1/3)A\,v_j,\qquad B\,v_j,
\]
and consequently
\[
 \operatorname{length}_R(M'_{P,j,+}/M'_{P,j,-})
       =\operatorname{length}_R(A/B)+r e_R,
 \qquad \eta'_{P,j}=\eta_P+r.                          \tag{G29.7}
\]
At a coefficient prime different from $3$, the same joint
character comparison has exactly the old length; the
algebra $T[C_3]$ is then the product of three copies of $T$
over a splitting coefficient ring, and each character factor
is treated separately.

The old length in (G29.7) can be arbitrarily large even in
the displayed local algebra fragment.  Indeed, for a
uniformizer $\pi$ and any integer $n\ge1$, take
\[
 T_n=M_n=\{(a,b)\in R^2:a\equiv b\pmod{\pi^n}\}.
\]
For the first-coordinate Hecke projector one has
\[
 A=R(1,0),\quad B=\pi^nR(1,0),\quad
       \eta_P=n/e_R,\quad \eta'_{P,j}=n/e_R+1.          \tag{G29.8}
\]
Thus the cubic splitting contributes a controlled additive
term and does not bound the original Hecke congruence length.
\end{proposition}
\begin{proof}
Since $R$ is Noetherian, $T\subset R^X$ is a finite
torsion-free $R$-module.  It is therefore finite free.
Point separation implies that $T\otimes_R\operatorname{Frac}R$
contains every coordinate idempotent.  For a fixed $x$, choose
for each $y\ne x$ an element $f_{xy}\in T$ with
$f_{xy}(x)\ne f_{xy}(y)$; the product of
\[
       \frac{f_{xy}-f_{xy}(y)}{f_{xy}(x)-f_{xy}(y)}
\]
over $y\ne x$ is that idempotent.  Thus the rational
eigenspace projectors used in the statement exist.

Evaluation of $\sum_{a=0}^2t_a\sigma^a$ at $(x,j)$
means $\sum_a t_a(x)\omega^{ja}$.  A polynomial of degree
at most two that vanishes at all three distinct roots is
zero over the fraction field.  Evaluation is consequently
injective.  Elements of $T$ separate unequal $x$'s, and
$\sigma$ separates unequal $j$'s, so $T'$ separates all
the stated characters.
For locality, every maximal ideal of the finite integral
extension $T'$ contracts to the unique maximal ideal
$\mathfrak m_T$ of $T$.
The residue field $k_T$ has characteristic $3$, and
\[
        T'/\mathfrak m_TT'=k_T[C_3]
                    \simeq k_T[z]/(z-1)^3
\]
has just one maximal ideal.  Hence $T'$ is local.
For faithfulness, apply $\sum_a t_a\sigma^a$ to
$m\otimes1$ for every $m\in M$.
The result is $\sum_a t_am\otimes\sigma^a$; its
vanishing implies that every $t_a$ acts as zero and hence
is zero.  This proves all asserted local algebra properties.

Identify $M'$ with three copies of $M$.
On one regular cubic orbit, $P_j$ has image
$(1/3)Rv_j$, with all three coordinates of $v_j$ units.
Thus applying the joint projector gives $(1/3)A v_j$.
A vector $a v_j$ in the joint rational eigenspace lies in
$M^3$ exactly when $a\in M$; since $a$ already belongs
to $\operatorname{im}P$, its intersection is $Bv_j$.
The exact sequence
\[
 0\longrightarrow A/B\longrightarrow (1/3)A/B
       \longrightarrow (1/3)A/A\longrightarrow0
\]
has last term isomorphic to $(R/3R)^r$.
Taking lengths and dividing by $e_R$ proves (G29.7).
If $3$ is a unit, the last term is zero and the usual
Chinese remainder theorem gives the three local character
factors as stated.

Finally $T_n$ is free with basis $(1,1)$ and $(\pi^n,0)$.
A pair in $T_n$ is a unit exactly when its common residue
modulo $\pi$ is nonzero: the coordinate inverses still
agree modulo $\pi^n$.  Hence $T_n$ is local.  The element
$(\pi^n,0)$ separates its two rational characters, and
its regular action on itself is faithful.  Projection of
$(1,1)$ gives $(1,0)$, while a pair $(a,0)$ lies in
$T_n$ exactly when $a\in\pi^nR$.
This proves (G29.8), and (G29.7) gives its joint cubic
version.  The example concerns precisely this local
algebra/lattice fragment; it is neither a countermodel to
all the fields of \texttt{NData} nor a refutation of a
supplied global theorem.
\end{proof}

\begin{proposition}[The full-packet regular tensor fails the existing good-row test]
\label{AC-G29-good-row}
Let $V$ be a nonzero representation of dimension $n$ over a
field, and let $P$ be the three-dimensional regular permutation
representation through a cubic character $\kappa$.
Then $V\otimes P$ is reducible: it has the canonical proper
invariant subspace
\[
             V\otimes\langle e_0+e_1+e_2\rangle.
                                                        \tag{G29.9}
\]
This remains true in characteristic $3$.
If $\kappa$ is trivial on a subgroup $I$, then for every
$\tau\in I$,
\[
 \det\bigl(X-\tau\mid V\otimes P\bigr)
                  =\det\bigl(X-\tau\mid V\bigr)^3.     \tag{G29.10}
\]
In particular this polynomial is not squarefree.

The source \texttt{CyclotomicNorms.StatementT.IsGood}
requires each prescribed auxiliary polynomial to be
squarefree.  A literal extension of an entire packet by
the regular cubic representation therefore cannot preserve
this good-row condition at an auxiliary inertia place
where $\kappa$ is unramified, if the prescribed polynomial
is transported through the existing inertia tie.
For $\chi_{269}$, every auxiliary rational prime different
from $269$ has this unramified property.
This rules out that particular full-packet substitution;
it does not rule out a new construction using separately
twisted packets or a regular action on a multiplicity space.
\end{proposition}
\begin{proof}
The vector $e_0+e_1+e_2$ is fixed by every permutation in
the regular action and is nonzero over every field.
Tensoring its one-dimensional span with $V$ gives an
invariant subspace of dimension $n$ in a space of
dimension $3n$, proving (G29.9).
If the action on $P$ is trivial on $I$, the restricted
representation is the direct sum of three copies of
$V|_I$.  Its characteristic polynomial is the cube in
(G29.10).  A nonconstant polynomial has a root over an
algebraic closure, and that root has multiplicity at least
three in its cube.  Thus the cube is not squarefree,
in every characteristic.

The source clause is explicit: following the monicity,
degree, and nonzero-constant-term conditions, \texttt{IsGood}
requires $\forall v,\ \texttt{Squarefree}(f(v))$.
The tie is also explicit in \texttt{DataN.NTypes}:
\texttt{inertiaCharpoly\_type} expresses the inertia
polynomial at the chosen tame generator as the product
of the linear factors obtained from the roots of
\texttt{items\_2}; \texttt{teichC\_inj} is injective on
the nonzero roots.  Consequently a squarefree monic
polynomial with nonzero constant term would give distinct
inertia eigenvalues under this tie.
The clause \texttt{NAut.wdOf\_ine} identifies those
inertia factors with the characteristic polynomial of an
actual Galois representation away from its coefficient
prime.  This is incompatible with (G29.10) at any such
auxiliary place.  Equivalently, a direct polynomial
transport would give $f(v)^3$, which fails the displayed
squarefree clause immediately.

The character $\chi_{269}$ is unramified away from $269$
by AC-K04 and AC-G22; restriction to a larger base field
preserves this unramified property.  The application to
the auxiliary inertia places follows.  The statement is
restricted to tensoring the entire Galois packet.  Applying
scalar twists to three separate packets leaves each
individual auxiliary inertia representation unchanged
and does not encounter this particular cubic repetition,
but requires the new cross-block arithmetic tie already
identified in AC-G29.
\end{proof}

\begin{remark}[The narrowed remaining theorem]
To use the regular principal construction in the supplied
arithmetic argument, one needs an actual mixed-block version
of its automorphic/cohomology dictionary: the new integral
row and theta lattices must be identified with the source
objects, the selected ranks must be controlled as in
(G28.6), and the strict local conditions and arithmetic tie
must be proved for the cross coefficients (G29.4).
The normalized extra character-projection loss can then be
bounded by AC-G28.  The uniform upper bound on those
cohomology lengths over the growing host family remains
a separate arithmetic assertion.  The construction does
not replace the field \texttt{Host.E} by an etale CM algebra
in an existing theorem, and its common scalar-twist variant
does not alter the pairwise pure-Tate coefficients at all.
\end{remark}
